\chapter{Of the Office of Mayor}
\label{art:mayor}

\articlenote{In which is set down the office of the Mayor, the civilian head of the silo, and the powers by which this office governs in concert with Judicial and Assembly.}

\section{The Office of Mayor}

The silo is governed by a civilian Mayor, who stands as the highest executive authority of the silo, answerable to the Assembly under Article 20 and to the requirements of this Pact as set down in Article 1. The Mayor is chosen by election and confirmed into office through the process set down in Article 5. The Mayor resides and keeps the office in the Up Top, in the Mayor's chambers, alongside the office of the Deputy Mayor.

\begin{clauses}
\item The Mayor is a citizen, subject to every duty this Pact lays upon any other citizen, save where an Article expressly grants the Mayor some further power or authority in the conduct of the office itself.
\item The Mayor holds office for a term of five years from confirmation, and continues in office until a successor is elected and confirmed under Article 5, an election being held as that term ends; a Mayor may stand again at the end of any term. Within a term the Mayor is removable by a vote of the Assembly as Article 20 provides.
\item In the absence or incapacity of the Mayor, the Deputy Mayor shall act as Mayor with all the powers of the office, for such time as the Mayor's absence continues, under Article 5.
\item The office of Mayor is barred to a citizen afflicted with the Syndrome, and vacated by it, as Article 13, Section 6 provides for every office of the silo; upon such a resignation the Deputy Mayor assumes the office under Article 5.
\end{clauses}

\section{The Mayor's Powers of Confirmation and Appointment}

The Mayor, in concert with this Pact and the Assembly, holds the power to confirm the heads of every Department and the single Judge who heads the Department of Judicial. The Sheriff is appointed by the Mayor alone, though the Mayor shall weigh the advice of those in and out of the Sheriff's own office. In all cases, the Mayor's judgment is final, and no other office may override such confirmation or appointment once given.

\begin{clauses}
\item The Mayor confirms the Head of each Department --- Mechanical, Information Technology, Supply, Mines, and all others established under this Pact --- by written order, countersigned, where a recommendation was made, by the officer who made it, and filed in the Judicial archive under Article 6.
\item The appointment of the Sheriff is the sole discretion of the Mayor: the outgoing Sheriff may make a recommendation, and the Mayor shall hear it in good faith, but the Mayor is not bound by that recommendation and may appoint any citizen of majority, whether from within the Sheriff's office or without, as the Mayor sees fit.
\item The confirmation of the Judge is made jointly by the Mayor and the Assembly under Article 20; the Judge is confirmed once in lifetime and stands until death or removal by the Assembly.
\item No Mayor shall dismiss or remove a confirmed head of Department from office for cause of ordinary policy disagreement. Removal for cause requires a finding of Judicial that the head has violated the Pact or a major duty of the office; such findings are referred to the Assembly under Article 20 for final removal.
\end{clauses}

\section{The Mayor's Orders and Emergency Authority}

The Mayor may issue standing orders, written and filed in the Judicial archive, directing the conduct of the silo's ordinary affairs in concert with the heads of the relevant Departments. No order of the Mayor may contradict this Pact, nor may any standing order be given in secret --- all orders are entered in the records that any citizen may examine under Article 6. In the event of emergency, the Mayor holds further powers, set down in Article 22.

\begin{clauses}
\item An order of the Mayor takes effect when issued and is binding on the head of the named Department; a Department head who disputes an order may refer the matter to Judicial for a ruling, and the order does not bind until Judicial has ruled.
\item The Mayor may not order any act that would require citizens to violate the Pact, nor direct the commission of any offense under Article 17. A citizen ordered to break the Pact may refuse the order and refer the matter to Judicial under Article 6.
\item The Mayor may make promises or take oath-bound obligations on behalf of the silo as a whole --- as, for example, the division of resources in a time of extreme shortage, or a standing obligation laid upon a Department beyond the term of the Mayor who laid it --- only with the consent of the Assembly, given under Article 20.
\item The Mayor directs the use of the silo's major reserves --- water, air, the generator's capacity, the stores of food --- in ordinary times as Article 11 provides, the Head of Supply setting the ration under the Mayor's general approval; in emergency, Article 22 governs.
\end{clauses}

\section{The Mayor and Judicial}

The Mayor and the Department of Judicial are co-equal in authority over the silo's law-making and law-executing functions, but with distinct roles: the Mayor commands the day-to-day operations and sets major policy; the Department of Judicial investigates offenses, tries accused citizens, and rules on the Pact's meaning when it is silent. Neither office may command the other in the exercise of this Pact-given power.

\begin{clauses}
\item The Mayor may refer a matter to Judicial for a ruling if the Pact is silent, and Article 6's procedure applies to it.
\item The Judge may, by written order, direct the Mayor to enforce a Judicial finding or to comply with a ruling made under Article 6; the Mayor must obey such a written order or refer the dispute to the Assembly under Article 20.
\item A dispute between the Mayor and Judicial over interpretation of the Pact is adjudicated by the Assembly under Article 20, which may convene an extraordinary session to resolve it.
\item The Mayor's orders do not bind Judicial's investigative or charging decisions; Judicial may investigate and charge any citizen, including the Mayor, for any offense under Article 17 without the Mayor's leave or permission.
\end{clauses}

\section{The Mayor and the Assembly}

The Mayor convenes the Assembly under Article 20 and sets its agenda, but does not control its votes or its business. The Mayor reports annually to the Assembly on the state of the silo. The Mayor may recommend amendments to this Pact under Article 23, and the Assembly may initiate amendments without the Mayor's consent.

\begin{clauses}
\item The Mayor is answerable to the Assembly for the conduct of the silo's affairs, and the Assembly may question the Mayor on any matter of policy or resource management.
\item No standing order, appointment, or decision of the Mayor binds the Assembly, and any citizen or the Assembly itself may petition for an extraordinary sitting to revisit a mayoral decision.
\end{clauses}

\section{The Office of Deputy Mayor and Succession}

The Deputy Mayor is the second-highest office of the silo, charged with assisting the Mayor in the conduct of the office and with assuming the mayoralty should the Mayor die, become incapacitated, or be removed by the Assembly. The office of Deputy Mayor, the manner of the Deputy's own coming to it, and the whole order of succession are set down in full in Article 5, which this Article's other provisions bind exactly as they bind the Mayor's own office.

\section{The Forgiveness Holiday}

The Mayor may declare a Forgiveness Holiday, of such duration and upon such notice as the Mayor sets, that a citizen holding a relic or restricted instrument or device forbidden under Article 16 may surrender it, in the manner that Article provides, without fear of answering for the mere keeping, making, or use of the thing surrendered.

\begin{clauses}
\item A citizen who surrenders a relic or restricted instrument or device under Article 16 during a Forgiveness Holiday is not answerable under Article 17 for the keeping, making, or use of the thing surrendered. The thing itself is examined and classified exactly as any other surrendered relic under Article 16, and a thing classified among the gravest is destroyed or sealed as that Article provides, whatever grace is given the citizen who surrendered it.
\item This grace covers only the keeping, making, or use of the thing surrendered, and no other offense the surrender, or Judicial's examination of the thing surrendered, may bring to light; a citizen who surrenders a relic is answerable exactly as any other citizen for any other violation of this Pact the surrender reveals.
\item A Forgiveness Holiday declared under this section is announced to the whole silo before it begins, that no citizen surrender in ignorance of the grace being offered, and its ending is announced with equal notice; a thing surrendered after a Forgiveness Holiday has ended is answered under Article 17 exactly as though no holiday had been declared.
\item The Mayor may declare a Forgiveness Holiday at such interval, and for such cause, as the Mayor judges the silo's peace to require, and need not state more than the fact of its declaration and its duration.
\end{clauses}
